
\documentclass{article}
\usepackage{colortbl}
\usepackage{makecell}
\usepackage{multirow}
\usepackage{supertabular}

\begin{document}

\newcounter{utterance}

\centering \large Interaction Transcript for game `bbh', experiment `ruin\_names', episode 26 with reflexion{-}llama{-}3.1{-}8b{-}instruct.
\vspace{24pt}

{ \footnotesize  \setcounter{utterance}{1}
\setlength{\tabcolsep}{0pt}
\begin{supertabular}{c@{$\;$}|p{.15\linewidth}@{}p{.15\linewidth}p{.15\linewidth}p{.15\linewidth}p{.15\linewidth}p{.15\linewidth}}
    \# & \multicolumn{2}{c}{Player} && \multicolumn{2}{c}{Game Master} \\
    \hline

    \theutterance \stepcounter{utterance}  
    & & & \multicolumn{4}{p{0.6\linewidth}}{
        \cellcolor[rgb]{0.9,0.9,0.9}{
            \makecell[{{p{\linewidth}}}]{
                \texttt{\tiny{[P1$\langle$GM]}}
                \texttt{Select the humorous edit that 'ruins' the input movie or musical artist name.} \\
\\ 
\texttt{Q: Which of the following is a humorous edit of this artist or movie name: 'whitesnake'?} \\
\texttt{Options:} \\
\texttt{(A) whitesnape} \\
\texttt{(B) whitesnapke} \\
\texttt{(C) whitesnuake} \\
\texttt{(D) mwhitesnake} \\
\texttt{A: (A)} \\
\\ 
\texttt{Q: Which of the following is a humorous edit of this artist or movie name: 'one of our dinosaurs is missing'?} \\
\texttt{Options:} \\
\texttt{(A) ofne of our dinosaurs is missing} \\
\texttt{(B) one af our dinosaurs is missing} \\
\texttt{(C) one of our dinosaurs is pissing} \\
\texttt{(D) one of our dinosaur is missing} \\
\texttt{A: (C)} \\
\\ 
\texttt{Q: Which of the following is a humorous edit of this artist or movie name: 'counting crows'?} \\
\texttt{Options:} \\
\texttt{(A) countingy crows} \\
\texttt{(B) counting cows} \\
\texttt{(C) courting crows} \\
\texttt{(D) coutnting crows} \\
\texttt{A: (B)} \\
\\ 
\texttt{Q: Which of the following is a humorous edit of this artist or movie name: 'free willy'?} \\
\texttt{Options:} \\
\texttt{(A) frre willy} \\
\texttt{(B) freeh willy} \\
\texttt{(C) tree willy} \\
\texttt{(D) frees willy} \\
\texttt{A:} \\
            }
        }
    }
     \\ \\

    \theutterance \stepcounter{utterance}  
    & \multicolumn{4}{p{0.6\linewidth}}{
        \cellcolor[rgb]{0.561,0.737,0.561}{
            \makecell[{{p{\linewidth}}}]{
                \texttt{\tiny{[P1$\rangle$GM]}}
                \texttt{Based on the pattern of humorous edits, the correct answer is (C) tree willy. This edit replaces the word "free" with "tree," creating a humorous and unexpected twist on the original name.} \\
            }
        }
    }
    & & \\ \\

    \theutterance \stepcounter{utterance}  
    & & & \multicolumn{2}{p{0.3\linewidth}}{
        \cellcolor[rgb]{0.9,0.9,0.9}{
            \makecell[{{p{\linewidth}}}]{
                \texttt{\tiny{[GM$|$GM]}}
                \texttt{Target: (C)} \\
            }
        }
    }
    & & \\ \\

    \theutterance \stepcounter{utterance}  
    & & & \multicolumn{2}{p{0.3\linewidth}}{
        \cellcolor[rgb]{0.9,0.9,0.9}{
            \makecell[{{p{\linewidth}}}]{
                \texttt{\tiny{[GM$|$GM]}}
                \texttt{game\_result = LOSE} \\
            }
        }
    }
    & & \\ \\

    \theutterance \stepcounter{utterance}  
    & & \multicolumn{4}{p{0.6\linewidth}}{
        \cellcolor[rgb]{0.95,0.95,0.95}{
            \makecell[{{p{\linewidth}}}]{
                \texttt{\tiny{[P1$\leftrightarrow$P]}}
                \texttt{Post{-}episode Reflexion (not sent to GM):} \\
\texttt{New plan: I failed to identify the correct humorous edit for the last question, 'free willy'. Upon reflection, I noticed that the correct humorous edits in the previous questions involved replacing a word with a word that is similar in sound or meaning, but with a twist. For example, 'whitesnape' replaces 'snake' with 'snape', which is a similar{-}sounding word. Similarly, 'one of our dinosaurs is pissing' replaces 'missing' with 'pissing', which is a humorous and unexpected twist. In the case of 'free willy', I should have looked for a similar{-}sounding word to replace 'free', such as 'tree' or 'fry'. However, I should have also considered other options, such as replacing 'free' with 'flee' or 'free' with 'fry'. Next time, I will make sure to consider multiple options and choose the one that is most likely to be the correct humorous edit.} \\
            }
        }
    }
    & \\ \\

\end{supertabular}
}

\end{document}
